\documentclass[10pt,a4paper]{amsart}
\usepackage[margin=19mm]{geometry}
\usepackage{amsmath,amssymb,mathtools}
\usepackage[scr=rsfs,cal=euler]{mathalfa}
\usepackage{xfrac}
\usepackage[unicode,hidelinks,bookmarks=false]{hyperref}
\newcommand{\ie}{i.e.\ }
\newcommand{\cf}{cf.\ }
\newcommand{\resp}{respectively\ }
\newcommand{\Subsection}{\subsection}
\providecommand{\nobreakdash}{\nobreak\hbox{-}}
\setlength{\emergencystretch}{2em}
\multlinegap0pt
\usepackage{amssymb}
\usepackage[shortlabels]{enumitem}
\usepackage{tikz}
\usepackage{bbm}
\usepackage{multirow}
\usepackage{xcolor}
\usepackage{booktabs}
\usepackage{algorithm}\usepackage{algpseudocode}
\usepackage{dsfont}
\usepackage{float}
\usepackage{caption}
\newtheorem{lemma}{Lemma}
\newcommand{\Q}{\widetilde{Q}}
\newcommand{\R}{\mathbb{R}}
\newcommand{\argmin}{\mathrm{argmin}  }
\newcommand{\ball}{\mathrm{B}}
\newcommand{\local}{\mathrm{loc}}
\newcommand{\norm}[1]{\left\lVert #1 \right\rVert}
\newcommand{\proj}{\operatorname{proj}}
\title{Constrained and Composite Sampling via Proximal Sampler}
\author{Thanh Dang, Jiaming Liang}
\date{Source: arXiv:2602.14478v1 (CC BY 4.0)}
\begin{document}
\maketitle

\noindent\emph{Source and licence.} Constrained and Composite Sampling via Proximal Sampler. Version arXiv:2602.14478v1 (CC BY 4.0), distributed by arXiv under the Creative Commons Attribution 4.0 International licence, http://creativecommons.org/licenses/by/4.0/, as recorded in the arXiv licence metadata for this exact version and shown on the abstract page https://arxiv.org/abs/2602.14478v1. Full licence notice: \texttt{licenses/arxiv-2602.14478-CC-BY-4.0.txt}.\par\smallskip

\noindent\textit{Editorial scope.} Counted unit: the article's lemma dolif:lem:localset (source0.tex lines 2721--2748). The article's complete original proof of it is delivered with it (source0.tex lines 2751--2863, one \texttt{proof} environment of 113 lines). No mathematics is written, repaired or re-derived: the statement and the proof are the article's own lines, in the article's own notation, and no retained line is substituted or re-flowed.  Retained uncounted as the source-local context the proof needs: lines 650--653; lines 672--676; lines 698--701; lines 2281--2288.  Nothing was omitted from any retained span, so the package carries no omission record.  No retained line contains a citation, so the wrapper carries no bibliography and no bibliography entry.  No retained line contains a figure environment, an image-inclusion command or drawing code, so no image is delivered and the package declares no asset dependency.  Editorial formatting only, in addition: an AMS \texttt{amsart} preamble that restates the article's own declarations and macros used by the retained lines, namely the environments \texttt{lemma} on separate counters, together with the standard packages those lines need.  The retained environments print in this document as Lemma 1 in order of appearance, whereas the article prints the same items with the numbering of its own full document.  The complete native source of the article as distributed in the arXiv e-print for this exact version is bundled unchanged as \texttt{sources/194711/source0.tex}.
\begin{equation}
\label{dolif:def:notation}
    p=(x,s,t), \, q=(y,u,v-b\eta);\quad \widetilde{C}_f=\sqrt{b^2+a^2+L_f^2}, \, \widetilde{C}_h=\sqrt{a^2+L_h^2}, \, \widetilde{C}_L=\widetilde{C}_f+\widetilde{C}_h. 
\end{equation}

\begin{equation}
\label{dolif:def:Theta}
    \widetilde{\Theta}^{\eta,\Q}_{y,u,v}(x,s,t)=I_{\Q}(x,s,t)+\frac{1}{2\eta}\|(x,s,t)-(y,u,v-\eta b)\|^2
+bv-\frac{\eta b^2}{2}. 
\end{equation}

\begin{equation}
\label{dolif:def:localball}
    \ball_{\local}:=2r_{\local}\ball_{d+2}(q), \qquad  \Q_{\local}:=\Q\cap \ball_{\local}.
\end{equation}

\begin{lemma}
\label{dolif:lem:minimizerinlocalball}  
Recall $\widetilde{\Theta}^{\eta,\Q}_{y,u,v}$ defined at~\eqref{dolif:def:Theta} and $\Q_{\local}=\Q\cap \ball_{\local}$ defined at~\eqref{dolif:def:localball}. The minimizer $ p_*=\argmin_{ p\in \R^{d+2}}\widetilde{\Theta}^{\eta,\Q}_{y,u,v}( p)=\proj_{\Q}( q)$
belongs to $\Q_{\local}$ a.s. In particular,
\[
 p_*=\underset{p\in \Q_\local}\argmin\widetilde{\Theta}^{\eta,\Q}_{y,u,v}( p).
\]
\end{lemma}

\begin{lemma}
\label{dolif:lem:localset}
Assume condition (B2) on $f$ and $h$. Fix any $c>1$ and define
\[
\ball_{c,\local}=c r_\local \ball_{d+2}( q),
\qquad
\Q_{\local}(c)=\Q\cap \ball_{c,\local},
\]
where $ q$ and $r_\local$ are defined at \eqref{dolif:def:localball}. Recall the notations $\widetilde{C}_h=\sqrt{L_h^2+a^2}, \widetilde{C}_f=\sqrt{L_f^2+a^2+b^2}$ at~\eqref{dolif:def:notation} and set
\[
\mu=\min\left\{\frac{a}{\widetilde{C}_h},\ \frac{b}{\widetilde{C}_f}\right\},\quad \lambda=2+\frac{a}{b},\quad \kappa=\frac{c-1}{\lambda+\mu}. 
\]
Moreover, define
\[
\bar p
=
 p_*+\left(0,\ \Delta s,\ \frac{a}{b}\Delta s+\Delta t\right),
\qquad
\Delta s=\Delta t=\kappa r_\local.
\]
Then we have the inclusion $B\left(\bar p,\ \kappa\mu r_\local\right)\subseteq \Q_{\local}(c)$. In particular,
\[
\mbox{\em minwidth}(\Q_{\local}(c))\ge 2\kappa\mu r_\local
\quad\text{and}\quad
\gamma(c):=\frac{c r_\local}{\mbox{\em minwidth}(\Q_{\local}(c))}
\le \frac{c(\lambda+\mu)}{(c-1)\mu}. 
\]
\end{lemma}

\begin{proof}  
Let us provide a sketch of the three steps of the proof.
To control $\gamma(c)=\frac{c r_{\local}}{\mbox{\em minwidth}(\Q_{\local}(c))}$, it suffices to lower bound
$\mbox{\em minwidth}(\Q_{\local}(c))$, which we do by constructing an inscribed ball inside
$\Q_{\local}(c)$.

Step~1 constructs a nontrivial ball $\ball_{\mathrm{in}}$ contained in $\Q$.
A subtlety is that Lemma~\ref{dolif:lem:minimizerinlocalball} only guarantees $ p_*\in \Q\cap r_{\local} \ball_{d+2}( q)$
and gives no quantitative slack for the constraints at $ p_*$; in particular, $ p_*$ may lie on or arbitrarily close to
$\partial\Q$, so a ball centered at $ p_*$ and inside $\Q$ need not exist.
We therefore shift $ p_*$ to
$\bar p= p_*+\left(0,\Delta s,\frac{a}{b}\Delta s+\Delta t\right)$ with $\Delta s=\Delta t=\kappa r_{\local}$,
which creates slack for both constraints: increasing $s$ relaxes $h(x)\le as$, while increasing $t$ by
$\frac{a}{b}\Delta s+\Delta t$ restores and adds slack for $f(x)+as\le bt$.
The Lipschitz bounds on $f$ and $h$ then imply that all points sufficiently close to $\bar p$ remain feasible,
yielding $\kappa\mu r_{\local} \ball_{d+2}(\bar p)\subseteq \Q$. Thus, $\kappa\mu r_{\local} \ball(\bar p)$ is the inscribed ball $\ball_{\mathrm{in}}$ in $\Q$ that we need. Note that this construction is nontrivial only when $c>1$, since then $\kappa=\frac{c-1}{\lambda+\mu}>0$.

Step~2 shows that $\ball_{\mathrm{in}}=\kappa\mu r_{\local}\ball_{d+2}(\bar p)$ also lies in $c r_{\local} \ball_{d+2}( q)$ and hence in $\Q_{\local}(c)$ (by combining with Step~1).
Using $\| p_*- q\|\le r_{\local}$ and
$\|\bar p- p_*\|=\Delta s+\left|\frac{a}{b}\Delta s+\Delta t\right|=\kappa\lambda r_{\local}$,
the triangle inequality gives
$\|\bar p- q\|+\kappa\mu r_{\local}\le r_{\local}+\kappa(\lambda+\mu)r_{\local}=c r_{\local}$. The last estimate implies $\kappa\mu r_{\local}\ball_{d+2}(\bar p)\subseteq c r_{\local}\ball_{d+2}( q)$. 

Step~3 uses the inclusion $\ball_{\mathrm{in}}=\kappa\mu r_{\local} \ball_{d+2}(\bar p)\subseteq \Q_{\local}(c)$ to lower bound $\mbox{\em minwidth}(\Q_{\local}(c))$ and finally upper bound $\gamma(c)$. 





\smallskip
\underline{Step 1:} Showing the nontrivial ball $\ball_{\mathrm{in}}:=\kappa\mu r_\local \ball(\bar p)\subseteq \Q$ for $c>1$.

\smallskip
Set
\[
g_1( p)=g_1(x,s,t)=h(x)-a s,
\qquad
g_2( p)=g_2(x,s,t)=f(x)+a s-b t,
\]
so that 
\begin{align}
\label{dolif:Qrewrite}
    \Q=\{ p:\ g_1( p)\le 0,\ g_2( p)\le 0\}. 
\end{align}
Because $h$ is $L_h$-Lipschitz and $f$ is $L_f$-Lipschitz, for all $ p=(x,s,t)$ and $ p'=(x',s',t')$,
\begin{align}
\label{dolif:consequencelipschitz}
|g_1( p)-g_1( p')|
\le L_h\|x-x'\|+a|s-s'|
\le \sqrt{L_h^2+a^2} \| p- p'\|
= \widetilde{C}_h\| p- p'\|,
\end{align}
and similarly $|g_2( p)-g_2( p')|\le \widetilde{C}_f\| p- p'\|$.

Since $ p_*\in \Q$, we have $g_1( p_*)\le 0$ and $g_2( p_*)\le 0$.
Define $\bar p$ as in the statement of the lemma then
\begin{align}
\label{dolif:compareg1g2}
g_1(\bar p)=g_1( p_*)-a\Delta s\le -a\Delta s,
\qquad
g_2(\bar p)=g_2( p_*)-b\Delta t\le -b\Delta t.
\end{align}
Moreover, we have 
\begin{align}
\label{dolif:compareradiusincribedball}
\kappa\mu r_\local=\min\left\{\frac{a\Delta s}{\widetilde{C}_h},\ \frac{b\Delta t}{\widetilde{C}_f}\right\}.
\end{align}
Then for any $\zeta$ with $\|\zeta\|\le \kappa\mu r_\local$, by \eqref{dolif:consequencelipschitz}, \eqref{dolif:compareg1g2}, \eqref{dolif:compareradiusincribedball}, we can write
\begin{align*}
g_1(\bar p+\zeta)&\stackrel{\eqref{dolif:consequencelipschitz}}{\le} g_1(\bar p)+\widetilde{C}_h\|\zeta\|\stackrel{\eqref{dolif:compareg1g2}}{\le} (-a\Delta s)+\widetilde{C}_h \kappa\mu r_\local\stackrel{\eqref{dolif:compareradiusincribedball}}{\le} 0,\\
g_2(\bar p+\zeta)&\le (-b\Delta t)+\widetilde{C}_f \kappa\mu r_\local\le 0. 
\end{align*}
Thus, in view of \eqref{dolif:Qrewrite}, we conclude $\kappa\mu r_\local \ball_{d+2}(\bar p)\subseteq \Q$. Notice in particular that the inscribed ball $\kappa\mu r_\local \ball_{d+2}(\bar p)$ is non-trivial only if $\kappa=\frac{c-1}{\lambda+\mu}>0$, and for that we require
\begin{align}
\label{dolif:whyc>1}
    c>1. 
    \end{align}

\smallskip
\underline{Step 2:} Showing $\ball_{\mathrm{in}}=\kappa\mu r_\local \ball_{d+2}(\bar p)\subseteq \Q_\local(c)$.

\smallskip

Next, since $\| p_*- q\|\le \| p_{k-1}- q\|=r_\local$, we have
\begin{align*}
\|\bar p- q\|
\le \| p_*- q\|+\|\bar p- p_*\|
&\le r_\local+\left(\Delta s+\left|\frac{a}{b}\Delta s+\Delta t\right|\right)\\
&= r_\local+\left(\left(1+\frac{a}{b}\right)\Delta s+\Delta t\right)
= r_\local+\kappa \lambda r_\local,
\end{align*}
which combined with $\kappa=\frac{c-1}{\lambda+\mu}$ leads to
\begin{align}
\label{dolif:inequalityconnectcenters}
\|\bar p- q\|+\kappa\mu r_\local
\le r_\local+\kappa(\lambda+\mu)r_\local
= r_\local+(c-1)r_\local
= c r_\local. 
\end{align}
Take any point $p\in \kappa\mu r_\local \ball_{d+2}(\bar p)$, then we get $\norm{p- q}\leq \norm{p-\bar p}+\norm{ q-\bar p}\leq  \kappa\mu r_\local+\|q-\bar p\|\stackrel{\eqref{dolif:inequalityconnectcenters}}{\leq} c r_\local$.
This implies $\kappa\mu r_\local \ball_{d+2}(\bar p)\subseteq c r_\local \ball_{d+2}( q)=\ball_{c,\local}$. Combining with Step~1 to get
$\kappa\mu r_\local \ball_{d+2}(\bar p)\subseteq \Q\cap \ball_{c,\local}=\Q_\local(c)$. 

\smallskip
\underline{Step 3:} Finally, since $\ball_{\mathrm{in}}=\kappa\mu r_\local \ball_{d+2}(\bar p)\subseteq \Q_\local(c)$ per Step 2, we have $\mbox{minwidth}(\Q_\local(c))\ge 2\kappa\mu r_\local$,
so that
\[
\gamma(c)=\frac{c r_\local}{\mbox{\em minwidth}(\Q_\local(c))}
\le \frac{c r_\local}{2\kappa\mu r_\local}
= \frac{c(\lambda+\mu)}{2(c-1)\mu}
\le \frac{c(\lambda+\mu)}{(c-1)\mu}.
\]
\end{proof}

\end{document}
